% !TEX program = pdflatex
% THEME 2 — Facile — Corrigé
\documentclass[11pt,a4paper]{article}
\usepackage[utf8]{inputenc}
\usepackage[T1]{fontenc}
\usepackage{lmodern}
\usepackage[french]{babel}
\usepackage{amsmath,amssymb,mathtools}
\usepackage{icomma}
\usepackage{siunitx}
\sisetup{output-decimal-marker={,},per-mode=symbol}
\DeclareSIUnit{\atm}{atm}
\DeclareSIUnit{\bar}{bar}
\DeclareSIUnit{\mmHg}{mmHg}
\usepackage[version=4]{mhchem}
\usepackage{xcolor}
\usepackage{colortbl}
\usepackage[most]{tcolorbox}
\usepackage{enumitem}
\usepackage{array}
\usepackage{booktabs}
\usepackage{fancyhdr}
\usepackage[a4paper,margin=2.1cm,headheight=26pt,headsep=10pt]{geometry}

\definecolor{primary}{RGB}{150,98,162}
\definecolor{primarydark}{RGB}{92,52,110}
\definecolor{excol}{RGB}{0,135,90}
\definecolor{attncol}{RGB}{200,40,55}

\newtcolorbox{rep}[1][]{enhanced,breakable,colback=excol!5,colframe=excol,
  boxrule=0.8pt,arc=2pt,left=3mm,right=3mm,top=2mm,bottom=2mm,
  fonttitle=\bfseries,coltitle=white,title={#1}}

\pagestyle{fancy}\fancyhf{}
\fancyhead[L]{\small\textcolor{primarydark}{\textbf{Fiche 7} — Loi des gaz}}
\fancyhead[R]{\small\textcolor{primarydark}{\textsc{Thème 2 · Facile · Corrigé}}}
\fancyfoot[C]{\small\thepage}
\renewcommand{\headrulewidth}{0.8pt}
\renewcommand{\headrule}{\hbox to\headwidth{\color{primary}\leaders\hrule height \headrulewidth\hfill}}
\setlength{\parindent}{0pt}\setlength{\parskip}{3pt}

\newcounter{exo}
\newenvironment{cor}[1][]{\refstepcounter{exo}\medskip
  \noindent{\color{primarydark}\bfseries Corrigé \theexo.}\ \textit{#1}\par\nopagebreak\smallskip}{\par}
\newcommand{\rf}[1]{\textbf{\textcolor{attncol}{#1}}}

\begin{document}

\begin{center}
{\Large\bfseries\color{primarydark} Thème 2 — Niveau Facile — Corrigé}
\end{center}
\textcolor{primary}{\hrule height 1pt}
\medskip

\begin{cor}[Combien de moles ?]
On cherche $n$ : on isole $n=\dfrac{pV}{RT}$. Unités atm/litre $\Rightarrow R=\num{0.082}$.
\[ n=\frac{pV}{RT}=\frac{\qty{2.0}{\atm}\times\qty{10.0}{\litre}}{\num{0.082}\times\qty{300}{\kelvin}}
    =\frac{20}{24{,}6}=\rf{\qty{0.81}{\mole}}. \]
\end{cor}

\begin{cor}[Quelle pression ?]
On isole $p=\dfrac{nRT}{V}$. Conversion : $T=27+273=\qty{300}{\kelvin}$.
\[ p=\frac{nRT}{V}=\frac{\qty{0.50}{\mole}\times\num{0.082}\times\qty{300}{\kelvin}}{\qty{5.0}{\litre}}
    =\frac{12{,}3}{5{,}0}=\rf{\qty{2.46}{\atm}}. \]
Soit environ deux fois et demie la pression atmosphérique.
\end{cor}

\begin{cor}[Quel volume ?]
On isole $V=\dfrac{nRT}{p}$.
\[ V=\frac{nRT}{p}=\frac{\qty{2.0}{\mole}\times\num{0.082}\times\qty{273}{\kelvin}}{\qty{1.0}{\atm}}
    =\rf{\qty{44.8}{\litre}}. \]
On reconnaît $2\times\qty{22.4}{\litre}$ : ce sont les conditions normales (voir Thème 3).
\end{cor}

\begin{cor}[Quelle température ?]
On isole $T=\dfrac{pV}{nR}$.
\[ T=\frac{pV}{nR}=\frac{\qty{1.0}{\atm}\times\qty{22.4}{\litre}}{\qty{1.0}{\mole}\times\num{0.082}}
    =\frac{22{,}4}{0{,}082}=\rf{\qty{273}{\kelvin}}, \]
soit $\theta=273-273=\rf{\qty{0}{\degreeCelsius}}$. C'est la définition même des conditions normales
(CNTP).
\end{cor}

\begin{cor}[Choisir le bon $R$]
\textbf{(a)} La pression est en \si{\kilo\pascal} et le volume en \si{\litre} : la valeur cohérente est
$R=\qty{8.314}{\kilo\pascal\litre\per\mole\per\kelvin}$. \emph{(Utiliser $\num{0.082}$ ici donnerait un
résultat faux d'un facteur $\sim100$.)}\\
\textbf{(b)}
\[ n=\frac{pV}{RT}=\frac{\qty{101.3}{\kilo\pascal}\times\qty{2.0}{\litre}}{\num{8.314}\times\qty{273}{\kelvin}}
   =\frac{202{,}6}{2269{,}7}=\rf{\qty{0.089}{\mole}}. \]
\end{cor}

\end{document}
